% Reader-facing Indonesian localization overlay for Notes on Diffy Qs v6.11.
% Load after diffyqssetup.sty.  This file changes presentation labels only;
% it does not change mathematical content, counters, labels, or references.

\renewcommand{\chaptername}{Bab}
\renewcommand{\appendixname}{Lampiran}
\renewcommand{\contentsname}{Daftar Isi}
\renewcommand{\listfigurename}{Daftar Gambar}
\renewcommand{\listtablename}{Daftar Tabel}
\renewcommand{\figurename}{Gambar}
\renewcommand{\tablename}{Tabel}
\renewcommand{\indexname}{Indeks}
\renewcommand{\bibname}{Daftar Pustaka}
\renewcommand{\partname}{Bagian}
\renewcommand{\proofname}{Bukti}
\renewcommand{\seename}{lihat}
\renewcommand{\alsoname}{lihat juga}

\renewcommand{\sectionnotes}[1]{%
  \noindent\emph{Catatan: #1}\medskip\par}

% amsthm creates these begin-environment commands directly.  Replacing only
% their displayed names preserves the theorem styles, counters, and optional
% notes established by diffyqssetup.sty.
\makeatletter
\renewcommand{\theorem}{%
  \@thm{\let\thm@swap\@gobble\th@plain}{theorem}{Teorema}}
\renewcommand{\exercise}{%
  \@thm{\let\thm@swap\@gobble\th@exercise}{exercise}{Latihan}}
\makeatother

\renewcommand{\exampleclaim}{%
  \par\addvspace{\medskipamount}\noindent\refstepcounter{example}%
  \hbox{\bfseries Contoh \theexample:}%
  \hspace{0.5em}\ignorespaces}
\renewcommand{\exampleclaimname}[1]{%
  \par\addvspace{\medskipamount}\noindent\refstepcounter{example}%
  \hbox{\bfseries Contoh \theexample\ {\normalfont (#1)}:}%
  \hspace{0.5em}\ignorespaces}
\renewcommand{\remarkclaim}{%
  \par\addvspace{\medskipamount}\noindent\refstepcounter{remark}%
  \hbox{\bfseries Catatan \theremark:}%
  \hspace{0.5em}\ignorespaces}
\renewcommand{\remarkclaimname}[1]{%
  \par\addvspace{\medskipamount}\noindent\refstepcounter{remark}%
  \hbox{\bfseries Catatan \theremark\ {\normalfont (#1)}:}%
  \hspace{0.5em}\ignorespaces}

\renewcommand{\figurevref}[1]{\myvref{Gambar}{#1}}
\renewcommand{\figureref}[1]{\myref{Gambar}{#1}}
\renewcommand{\tablevref}[1]{\myvref{Tabel}{#1}}
\renewcommand{\tableref}[1]{\myref{Tabel}{#1}}
\renewcommand{\chapterref}[1]{\myref{bab}{#1}}
\renewcommand{\Chapterref}[1]{\myref{Bab}{#1}}
\renewcommand{\appendixref}[1]{\myref{lampiran}{#1}}
\renewcommand{\Appendixref}[1]{\myref{Lampiran}{#1}}
\renewcommand{\subsectionref}[1]{\myref{subbagian}{#1}}
\renewcommand{\subsectionvref}[1]{\myvref{subbagian}{#1}}
\renewcommand{\exercisevref}[1]{\myvref{Latihan}{#1}}
\renewcommand{\exerciseref}[1]{\myref{Latihan}{#1}}
\renewcommand{\examplevref}[1]{\myvref{Contoh}{#1}}
\renewcommand{\exampleref}[1]{\myref{Contoh}{#1}}
\renewcommand{\thmvref}[1]{\myvref{Teorema}{#1}}
\renewcommand{\thmref}[1]{\myref{Teorema}{#1}}
\renewcommand{\remarkref}[1]{\myref{Catatan}{#1}}

% varioref page-position phrases used by \myvref.
\renewcommand{\reftextfaceafter}{pada halaman seberang}
\renewcommand{\reftextfacebefore}{pada halaman seberang}
\renewcommand{\reftextafter}{pada halaman berikutnya}
\renewcommand{\reftextbefore}{pada halaman sebelumnya}
\renewcommand{\reftextcurrent}{pada halaman ini}
\renewcommand{\reftextfaraway}[1]{pada halaman~\pageref{#1}}

\renewcommand{\printanswers}{%
  \cleardoublepage
  \phantomsection
  \chapter*{Solusi untuk Latihan Terpilih}%
  \markboth{SOLUSI UNTUK LATIHAN TERPILIH}{SOLUSI UNTUK LATIHAN TERPILIH}%
  \addfakecontentsline{Solusi untuk Latihan Terpilih}%
  \unvbox\answerscollect\setbox\answerscollect=\vbox{}% 
}
